\section{Conclusion}

This paper broadens the applicability of equality saturation by making it scale to complex optimizations in a functional language in two ways.

Firstly, sketch-guided equality saturation is proposed as a semi-automated technique for scaling to more complex optimizations by factoring a single equality saturation search into a sequence of smaller equality saturation searches (\cref{sketching}).
Programmers guide the process by describing how a program should evolve using a sequence of sketches.
We demonstrate that sketch-guided equality saturation enables seven complex optimizations of matrix multiplication to be applied within seconds in the \Rise{} functional language, using under 1 GB of RAM, using no more than three sketch guides (\cref{fig:s3}).
By contrast, traditional unguided equality saturation cannot discover the five more complex optimizations even with an hour of runtime and 60 GB of RAM (\cref{fig:s2}).
For each optimization, the generated code is identical to the high-performance code generated by manually ordering thousands of rewrites via rewriting strategies.
The performance of this code is comparable with code produced by the state-of-the-art TVM compiler~\cite{hagedorn2020-elevate}.
% While complex optimizations of matrix multiplication are not feasible with unguided equality saturation, they are feasible with sketch-guided equality saturation
% Equality saturation alone does not scale to these optimizations, even with an hour of compilation time and 60GB of RAM.

Secondly, to effectively apply equality saturation to functional languages, we propose new techniques to efficiently encode lambda calculi for equality saturation. 
The key innovations are extraction-based substitution and representing identifiers as De Bruijn indices (\cref{eqsat-bindings}). Specifically for our  \Rise{} use case, we apply the techniques to a polymorphically typed lambda calculus. Combining the techniques reduces the runtime and memory consumption of equality saturation over lambda terms by orders of magnitude, enabling equality saturation to scale to more complex optimization goals  (\cref{fig:lambda-res}).

\medskip

Future work may investigate how to identify appropriate sets of rules to be applied in each search.
We are interested in exploring how to design effective sketch guides for more diverse applications, and to explore ways to possibly even synthesize sketch guides automatically.
% \pwtcomment{There are too many ideas in the preceding sentence. Split into multiple sentences}
%We believe that rewriting strategies would benefit from being more declarative, for example via sketch guidance, as opposed to being procedural step-by-step recipes.
%Furthermore, it would be interesting to integrate sketch guidance into an interactive optimization assistant.
We believe that combining imperative, step-by-step approaches (e.g. rewriting strategies) with more declarative approaches (e.g. sketch guidance) deserves further research and can lead to the creation of practical, interactive optimization assistants.
Finally, we hope that the community will be inspired to apply and extend sketch-guided equality saturation in new domains.
% We hope that this paper will encourage the community to create variations of sketch-guided equality saturation and of the \sLang{} language for their own use cases.
% We hope that the paper will spark future work on more diverse applications.

% Future Work: combine with prioritizing e-graph growth; study how to come up with good sets of rules / sketch guides; evaluate usefuleness in other settings than Rise opt; extend SketchBasic (e.g. 'and' sketch); integrate SGES with ELEVATE to allow even more diverse trade-offs; synthesize sketch guides to only require sketch goal; study impact of loss of completeness in encoding}
%  In the future, we envision using strategies to offer tradeoffs between precise control (as in this paper) and full automation (as in LIFT). Users would then provide strategies that constrain and guide the automated exploration process.

\clearpage